\chapter{Details of Preconditioning and Statistical Analysis}
\label{app:reconstruction_details}

\section{Assembly of the Lewis--Sendov Voxel Blocks}
\label{app:ls_block_assembly}

The block approximation used in Chapter~\ref{chapter:optimisation} retains PET/SPECT coupling within each voxel and omits coupling between voxels. The following construction specifies how the local derivatives are assembled into those approximate blocks.

In practice each voxel block is assembled per neighbourhood target: the Lewis--Sendov calculus is evaluated in local Jacobian coordinates, and the retained direction-diagonal blocks are mapped into modality space by the diagonal congruence
\begin{equation}
    T_{r,\ell}
    =
    \operatorname{diag}(a_{r,\ell})\,
    \Theta_{r,\ell}\,
    \operatorname{diag}(a_{r,\ell})
    =
    \bigl(a_{r,\ell}a_{r,\ell}^{\top}\bigr)\circ \Theta_{r,\ell},
    \label{eq:ls_block_pullback}
\end{equation}
where $\Theta_{r,\ell}$ is the local Hessian block for difference channel $\ell$ at target $r$ and the modality scalings $[a_{r,\ell}]_m$ combine the modality weights $\rho_{m,r}$ and the stencil step scaling with the square roots of the channel-$\ell$ diagonal entries of the operator Gram matrices $\Pi_{m,r}^{\top}\Pi_{m,r}$. Each contribution~\eqref{eq:ls_block_pullback} is the Hadamard product of a rank-one Gram matrix with a positive-semidefinite local block---$\Theta_{r,\ell}$ is a principal block of the local Hessian of the convex smoothed spectral penalty---hence positive semidefinite by the Schur product theorem. Accumulated over the targets incident on voxel $n$ and their difference channels, the assembled voxel blocks are therefore positive semidefinite before flooring. Because the scalings compress the cross-operator geometry to its diagonal energies, the assembled blocks approximate the exact voxel blocks $[H_{\mathcal R}]_{nn}$ rather than reproducing them exactly.

\section{Statistical Models for the Phantom and Clinical Comparisons}
\label{app:90y_statistical_models}

These are the full model specifications for the comparisons in Chapter~\ref{chapter:paper}. In the phantom model, $c$ identifies the reconstruction method, $v$ the sphere, and $\mathrm{bs}$ the bootstrap replicate. In the clinical models, $p$ identifies the patient and $l$ the lesion. The reference method is \wdtnv. The symbols $\beta_0$ and $\alpha_c$ denote the intercept and method effects within each model; parameters are estimated separately for each endpoint. Random effects and residual errors are assumed mutually independent, with the Gaussian distributions specified below.

\subsection{Phantom Recovery Coefficients}

Inference for RC uses a linear mixed model on
\begin{equation}
y_{\mathrm{bs},c,v} := \log\!\left(\mathrm{RC}^{(c)}_{\mathrm{bs},v}\right),
\qquad \mathrm{bs}=1,\ldots,N_{\mathrm{bs}},
\end{equation}
with a random intercept for bootstrap replicate and fixed effects for method, sphere, and their interaction:
\begin{align}
y_{\mathrm{bs},c,v} &= \beta_0 + \alpha_c + \gamma_v + (\alpha\gamma)_{c,v} + u_{\mathrm{bs}} + \varepsilon_{\mathrm{bs},c,v}, \nonumber \\
u_{\mathrm{bs}} &\overset{\mathrm{iid}}{\sim} \mathcal{N}(0,\tau_{\mathrm{bs}}^2),
\qquad
\varepsilon_{\mathrm{bs},c,v} \overset{\mathrm{iid}}{\sim} \mathcal{N}(0,\tau^2).
\end{align}

\subsection{Clinical Contrast and Background Noise}

We analyse $\log(\mathrm{TBR})$ using a linear mixed-effects model with fixed method effects and random intercepts for patient and lesion-within-patient:
\begin{equation}
\begin{aligned}
y_{p,l,c} &:= \log\!\left(\mathrm{TBR}_{p,l,c}\right)
= \beta_0 + \alpha_c + u_p + v_{p,l} + \varepsilon_{p,l,c},\\
u_p &\overset{\mathrm{iid}}{\sim} \mathcal{N}(0,\tau^2_{\mathrm{p}}),
\qquad
v_{p,l} \overset{\mathrm{iid}}{\sim} \mathcal{N}(0,\tau^2_{\mathrm{l}}),
\qquad
\varepsilon_{p,l,c} \overset{\mathrm{iid}}{\sim} \mathcal{N}(0,\tau^2).
\end{aligned}
\end{equation}
Background noise is analysed analogously at the patient level using
\begin{equation}
\begin{aligned}
y^{\mathrm{CoV}}_{p,c} &:= \log\!\left(\mathrm{CoV}_{p,c}(\%)\right)
= \beta_0 + \alpha_c + u_p + \varepsilon_{p,c},\\
u_p &\overset{\mathrm{iid}}{\sim} \mathcal{N}(0,\tau^2_{\mathrm{p,CoV}}),
\qquad
\varepsilon_{p,c} \overset{\mathrm{iid}}{\sim} \mathcal{N}(0,\psi^2).
\end{aligned}
\end{equation}
